% !TeX root = ../../Book4.tex
% 纲要标题：### 3.4 伴随的后果
% 纲要位置：计划纲要.md，第 3177 行。

\section{伴随的后果}
\label{b4:sec:0304}

伴随双射可以直接改写极限或余极限的泛性质，从而确定哪些构造被保留。单位与余单位的复合还产生单子结构，说明连续施加伴随函子时出现的代数规律。

% 正文按纲要写入；各节的基础结果在本章证明。

% 纲要内容（仅作写作计划，尚非教材正文）：
%
% * 左伴随保持余极限
% * 右伴随保持极限
% * 一般伴随函子定理的极限保持、解集和选择条件，完整证明置于附录
% * 单子与余单子初步，完整单子性理论置于附录
% ---
%

\subsection{左伴随与余极限}
\label{b4:subsec:0304:01}

余极限的泛性质通过从所得对象出发的映射表述，伴随双射恰好能改写这种映射。下面的结论因此不依赖余极限在某个范畴中的具体构造。

\begin{theorem}\label{b4:thm:adjoints-preserve-colimits}
设 \(\mathcal C,\mathcal D\) 为局部小范畴，\(F:\mathcal C\rightarrow\mathcal D\) 有右伴随 \(G:\mathcal D\rightarrow\mathcal C\)。若小范畴 \(J\) 上的图表 \(K:J\rightarrow\mathcal C\) 有余极限 \((L,(\iota_j)_{j\in J})\)，则 \((FL,(F\iota_j)_{j\in J})\) 是 \(FK\) 的余极限。因此左伴随保持所有存在的小余极限，包括空图表的余极限。
\end{theorem}
\begin{proof}
记伴随在两变量自然的双射为
\[
\Phi_{X,Y}:\operatorname{Hom}_{\mathcal D}(FX,Y)
\xrightarrow{\sim}\operatorname{Hom}_{\mathcal C}(X,GY),
\qquad X\in\mathcal C,\ Y\in\mathcal D.
\]
函子保持复合，故 \((F\iota_j)\) 是余锥。要验证它的泛性质，可以依次改写从 \(FL\) 出发的 Hom 集：
\[
\begin{aligned}
\operatorname{Hom}_{\mathcal D}(FL,Y)
&\cong\operatorname{Hom}_{\mathcal C}(L,GY)\\
&\cong\{\,K\;\text{到}\;GY\;\text{的余锥}\,\}\\
&\cong\{\,FK\;\text{到}\;Y\;\text{的余锥}\,\}.
\end{aligned}
\]
第一步是伴随双射，第二步是 \(L\) 的余极限泛性质，第三步逐分量使用伴随双射；自然性保证余锥的相容条件也随之对应。这些双射关于 \(Y\) 自然。具体地，取余锥 \(q_j:FK(j)\rightarrow Y\)，令
\[
s_j=\Phi_{K(j),Y}(q_j):K(j)\longrightarrow GY.
\]
对 \(a:i\rightarrow j\)，余锥相容性为 \(q_jFK(a)=q_i\)。伴随在第一变量的自然性给出 \(s_jK(a)=s_i\)，所以 \((s_j)\) 是 \(K\) 的余锥。由 \(L\) 的泛性质，存在唯一 \(s:L\rightarrow GY\) 满足 \(s\iota_j=s_j\)。令 \(q=\Phi^{-1}_{L,Y}(s):FL\rightarrow Y\)。再次利用自然性，
\[
\Phi_{K(j),Y}\bigl(qF(\iota_j)\bigr)
=\Phi_{L,Y}(q)\iota_j=s\iota_j=s_j=\Phi_{K(j),Y}(q_j),
\]
双射性给出 \(qF(\iota_j)=q_j\)。若 \(q'\) 也满足这些等式，则 \(\Phi(q')\) 也经过每个 \(\iota_j\) 给出 \(s_j\)，由余极限的唯一性有 \(\Phi(q')=s\)，从而 \(q'=q\)。证明没有要求 \(J\) 非空；在空图表情形，它就是始对象的泛性质。
\end{proof}

例如 \({}_SP_R\) 固定时，\cref{b4:prop:hom-tensor-adjunction}表明 \(P\otimes_R-\) 保持任意直和、余等化子及滤过余极限。这并非说所有相关对象的底集合都按同一公式构造，而是说它们具有同一个余极限泛性质。具体的比较同构由结构映射唯一确定，例如
\[
\bigoplus_i(P\otimes_RM_i)\xrightarrow{\sim}P\otimes_R\left(\bigoplus_iM_i\right)
\]
在第 \(i\) 个直和项上把 \(p\otimes m\) 送到 \(p\otimes\iota_i(m)\)。

\ProseParagraphBreak
左伴随的条件本身不保证保持积。\cref{b4:exm:exact-not-products}已经用共同分母证明，\(\mathbb Q\otimes_{\mathbb Z}-:\mathbf{Ab}\rightarrow\mathbb Q\text{-}\mathbf{Mod}\) 不保持可数个 \(\mathbb Z\) 的积。现在将\cref{b4:prop:hom-tensor-adjunction}用于 \({}_{\mathbb Q}\mathbb Q_{\mathbb Z}\)，又知道它是左伴随，因此这个例子也检验了保持定理的方向：它保持余极限，却仍然可能不保持无限积。此处不能从无限积失败推断有限积也失败。

\subsection{右伴随与极限}
\label{b4:subsec:0304:02}

\begin{theorem}\label{b4:thm:adjoints-preserve-limits}
设 \(\mathcal C,\mathcal D\) 为局部小范畴，\(F:\mathcal C\rightarrow\mathcal D\) 左伴随于 \(G:\mathcal D\rightarrow\mathcal C\)。若小图表 \(K:J\rightarrow\mathcal D\) 有极限 \((L,(p_j)_{j\in J})\)，则 \((GL,(Gp_j)_{j\in J})\) 是 \(GK\) 的极限。因此右伴随保持所有存在的小极限，包括空图表的极限。
\end{theorem}
\begin{proof}
记伴随在两变量自然的双射为
\[
\Phi_{X,Y}:\operatorname{Hom}_{\mathcal D}(FX,Y)
\xrightarrow{\sim}\operatorname{Hom}_{\mathcal C}(X,GY),
\qquad X\in\mathcal C,\ Y\in\mathcal D.
\]
这次改写的是指向极限的 Hom 集：
\[
\begin{aligned}
\operatorname{Hom}_{\mathcal C}(X,GL)
&\cong\operatorname{Hom}_{\mathcal D}(FX,L)\\
&\cong\{\,FX\;\text{到}\;K\;\text{的锥}\,\}\\
&\cong\{\,X\;\text{到}\;GK\;\text{的锥}\,\}.
\end{aligned}
\]
中间一步使用 \(L\) 的极限泛性质，两侧使用伴随双射，所有对应关于 \(X\) 自然。具体地，给定锥 \(s_j:X\rightarrow GK(j)\)，用逆伴随双射得到 \(q_j:FX\rightarrow K(j)\)。对 \(a:i\rightarrow j\)，关系 \(GK(a)s_i=s_j\) 通过第二变量自然性等价于 \(K(a)q_i=q_j\)。因此 \((q_j)\) 为锥，唯一确定 \(q:FX\rightarrow L\)，使 \(p_jq=q_j\)。令 \(s=\Phi_{X,L}(q)\)，第二变量自然性给出
\[
Gp_j\circ s=\Phi_{X,K(j)}(p_jq)=s_j.
\]
任何具有这些复合的 \(s'\)，在逆双射下都给出具有投影 \(q_j\) 的态射 \(q'\)，由 \(L\) 的唯一性有 \(q'=q\)，进而 \(s'=s\)。空图表情形同样成立，给出终对象的保持。
\end{proof}

固定源的 Hom 函子保持极限已在\cref{b4:prop:representables-preserve-limits}直接证明：对于有极限的图表 \(K:J\rightarrow\mathcal C\) 及 \(A\in\mathcal C\)，一个态射 \(A\rightarrow\lim_JK\) 恰好对应从 \(A\) 到各 \(K(j)\) 的相容态射族，所以
\[
\operatorname{Hom}_{\mathcal C}(A,\lim_JK)
\cong\lim_{j\in J}\operatorname{Hom}_{\mathcal C}(A,K(j)).
\]
若一个函子与固定源的 Hom 函子自然同构，相应的锥经分量同构转换，仍具有相同的唯一分解性质，所以任意协变可表函子保持极限。反变可表函子则把原范畴的余极限变成集合的极限；必须连同相反范畴一起读这个断言，不能遗漏方向。

\ProseParagraphBreak
由 Hom–张量伴随，\(\operatorname{Hom}_S(P,-)\) 保持积与等化子。若给定左 \(S\) 模同态 \(u:N\rightarrow N'\)，它的核是 \(u\) 与零映射的等化子，故
\[
\operatorname{Hom}_S(P,\ker u)
\cong\ker\bigl(\operatorname{Hom}_S(P,u)\bigr).
\]
右边由那些满足 \(u\circ f=0\) 的 \(f:P\rightarrow N\) 组成。一般情形下，\(\operatorname{Hom}_S(P,-)\) 未必保持满同态；例如 \(\mathbb Z\rightarrow\mathbb Z/n\mathbb Z\) 在 \(n>1\) 时满射，而
\[
\operatorname{Hom}_{\mathbb Z}(\mathbb Z/n\mathbb Z,\mathbb Z)=0,
\qquad
\operatorname{Hom}_{\mathbb Z}(\mathbb Z/n\mathbb Z,\mathbb Z/n\mathbb Z)
\cong\mathbb Z/n\mathbb Z.
\]
因此诱导映射不是满射。保持核与有限积使这个 Hom 函子具有\secref{b4:sec:0204}中的左正合性；上例的满射端失败，说明它一般不右正合。下一章会把这些正合性质放到一般交换范畴中讨论。

\ProseParagraphBreak
保持极限是一项必要条件，并非任意情况下都足以推出存在左伴随。证明存在伴随还要为每个对象构造相应的表示对象，并控制候选对象的大小。设 \(\mathcal C,\mathcal D\) 为局部小范畴，\(\mathcal D\) 有全部小极限，\(U:\mathcal D\to\mathcal C\) 为函子。解集条件要求：对每个 \(X\in\mathcal C\)，存在集合大小的一族箭头 \(u_i:X\to UY_i\)，使每个 \(u:X\to UY\) 都能写成 \(u=U(v)u_i\)，其中 \(v:Y_i\to Y\)。它不要求这批候选已经具有唯一分解性质，而是从可能遍历真类的候选中选出集合大小的一族，足以产生全部这些箭头。再假定满足解集条件时，所需的解集与小极限可以按对象类同时选择。\Cref{b4:thm:A-general-adjoint-functor}将在这些条件下证明：\(U\) 有左伴随，当且仅当它保持全部小极限并满足解集条件。

\subsection{单子与余单子初步}
\label{b4:subsec:0304:03}

对伴随 \(F\dashv G\)，复合 \(T=GF\) 是 \(\mathcal C\) 上的自函子；\(T^2\) 表示函子复合 \(T\circ T=GFGF\)。单位给出 \(X\rightarrow TX\)。余单位在 \(FX\) 处的分量 \(\varepsilon_{FX}:FGFX\rightarrow FX\) 将中间的 \(FG\) 归并，施加 \(G\) 后得到
\[
\mu_X=G(\varepsilon_{FX}):T^2X=GFGFX\longrightarrow GFX=TX.
\]
在自由左 \(R\) 模与底集合函子的伴随中，\(TX\) 是 \(X\) 的形式有限线性组合，而 \(T^2X\) 是“形式线性组合的形式线性组合”。\(\mu_X\) 展开括号并合并同类项。这提示我们暂时只保留自函子、插入元素和归并运算。

\begin{definition}\label{b4:def:monad}
设 \(\mathcal C\) 为范畴，\(T:\mathcal C\rightarrow\mathcal C\) 为自函子，\(\eta:1_{\mathcal C}\Rightarrow T\)、\(\mu:T^2\Rightarrow T\) 为自然变换。若对每个 \(X\in\mathcal C\) 都有
\[
\mu_X\circ T(\mu_X)=\mu_X\circ\mu_{TX},
\qquad
\mu_X\circ T(\eta_X)=1_{TX}=\mu_X\circ\eta_{TX},
\]
则称 \((T,\eta,\mu)\) 为 \(\mathcal C\) 上的\term[danzi]{单子}[monad]。第一式称结合律，后两式称单位律。
\end{definition}
\symidx{\(\mu:T^2\Rightarrow T\)：单子的乘法自然变换}

\[
\begin{tikzcd}[column sep=large,row sep=large]
T^3X\arrow[r,"T(\mu_X)"]\arrow[d,"\mu_{TX}"']&
T^2X\arrow[d,"\mu_X"]\\
T^2X\arrow[r,"\mu_X"']&TX
\end{tikzcd}
\]
\[
\begin{tikzcd}[column sep=large,row sep=large]
TX\arrow[r,shift left=0.7ex,"T(\eta_X)"]
\arrow[r,shift right=0.7ex,"\eta_{TX}"']
\arrow[dr,"1_{TX}"']&T^2X\arrow[d,"\mu_X"]\\
&TX.
\end{tikzcd}
\]

结合律的方块比较从三层复合到一层复合的两种归并次序。单位律的图则有两个不同的插入位置：\(T(\eta_X)\) 在已经施加的 \(T\) 内部插入单位，\(\eta_{TX}\) 对整个 \(TX\) 插入单位。二者有相同的源与目标，却未必是同一态射；单位律要求它们分别与 \(\mu_X\) 复合后都还原 \(TX\)。

\begin{proposition}\label{b4:prop:adjunction-monad}
设 \(F:\mathcal C\rightarrow\mathcal D\)、\(G:\mathcal D\rightarrow\mathcal C\) 为局部小范畴之间的伴随，其单位、余单位为 \(\eta,\varepsilon\)。则
\[
T=GF,\qquad \mu_X=G(\varepsilon_{FX})
\]
与原有单位 \(\eta\) 组成 \(\mathcal C\) 上的单子。
\end{proposition}
\begin{proof}
对 \(f:X\rightarrow X'\)，要证明 \(T(f)\mu_X=\mu_{X'}T^2(f)\)。将 \(F(f):FX\rightarrow FX'\) 代入余单位的自然性，得到
\[
F(f)\circ\varepsilon_{FX}
=\varepsilon_{FX'}\circ FG(F(f)).
\]
施加 \(G\)，左边为 \(T(f)\mu_X\)，右边为 \(\mu_{X'}T^2(f)\)，故 \(\mu\) 自然。两条单位律分别为
\[
\begin{aligned}
\mu_XT(\eta_X)
&=G\bigl(\varepsilon_{FX}F(\eta_X)\bigr)=1_{GFX},\\
\mu_X\eta_{TX}
&=G(\varepsilon_{FX})\eta_{GFX}=1_{GFX},
\end{aligned}
\]
第一行对 \(\varepsilon_{FX}F(\eta_X)=1_{FX}\) 施加 \(G\)，第二行将另一个三角恒等式中的 \(Y\) 取为 \(FX\)。

结合律也来自余单位的自然性，但这次使用的态射是
\(u=\varepsilon_{FX}:FGFX\rightarrow FX\)。对这两个对象及 \(u\)，自然性给出
\[
u\circ\varepsilon_{FGFX}
=\varepsilon_{FX}\circ FG(u).
\]
施加 \(G\) 后，\(G(u)=\mu_X\)，而
\[
G(\varepsilon_{FGFX})=\mu_{TX},\qquad
G(FG(u))=T(\mu_X).
\]
因此得到 \(\mu_X\mu_{TX}=\mu_XT(\mu_X)\)，正是结合律。
\end{proof}

在自由左 \(R\) 模伴随中，设 \(R\) 为任意含幺环。令 \(e_x\) 表示 \(TX\) 中的形式生成元，而 \(e_{v_i}\) 表示 \(T^2X\) 中由 \(v_i\in TX\) 标记的生成元，则
\[
\mu_X\left(\sum_i r_i e_{v_i}\right)
=\sum_i r_iv_i
=\sum_x\left(\sum_i r_is_{i,x}\right)e_x,
\qquad v_i=\sum_xs_{i,x}e_x.
\]
所有和都有限，标量次序为 \(r_is_{i,x}\)，不能在非交换环中调换。对 \(v=\sum_xs_xe_x\in TX\)，两种插入方式具体为
\[
T(\eta_X)(v)=\sum_xs_xe_{e_x},\qquad
\eta_{TX}(v)=e_v.
\]
前者把每个原生成元再包成一项，后者把整个线性组合当成一个新生成元；\(\mu_X\) 分别作用后都得到 \(v\)。

\ProseParagraphBreak
三层的结合律也可以直接计算。取 \(X=\{x,y,z\}\) 及 \(r,s,t\in R\)，令
\[
v=s e_x+e_y\in TX,\qquad
w=r e_v+e_{e_z}\in T^2X,\qquad
\xi=t e_w\in T^3X.
\]
按两种次序归并，分别得到
\[
\begin{aligned}
\mu_XT(\mu_X)(\xi)
&=\mu_X\bigl(t e_{rs e_x+r e_y+e_z}\bigr)\\
&=t(rs)e_x+(tr)e_y+t e_z,\\
\mu_X\mu_{TX}(\xi)
&=\mu_X\bigl((tr)e_v+t e_{e_z}\bigr)\\
&=((tr)s)e_x+(tr)e_y+t e_z.
\end{aligned}
\]
环的结合律给出 \(t(rs)=(tr)s\)，所以两条路线的结果相同，且整个计算没有交换任何标量的次序。

\begin{definition}\label{b4:def:comonad}
设 \(\mathcal D\) 为范畴，\(Q:\mathcal D\rightarrow\mathcal D\) 为自函子，\(\varepsilon:Q\Rightarrow1_{\mathcal D}\)、\(\delta:Q\Rightarrow Q^2\) 为自然变换。若对每个 \(Y\in\mathcal D\) 都有
\[
Q(\delta_Y)\circ\delta_Y=\delta_{QY}\circ\delta_Y,
\qquad
Q(\varepsilon_Y)\circ\delta_Y=1_{QY}
=\varepsilon_{QY}\circ\delta_Y,
\]
则称 \((Q,\varepsilon,\delta)\) 为 \(\mathcal D\) 上的\term[yudanzi]{余单子}[comonad]。
\end{definition}
\symidx{\(\delta:Q\Rightarrow Q^2\)：余单子的余乘法自然变换}

伴随 \(F\dashv G\) 在另一范畴上给出 \(Q=FG\)，余单位仍为 \(\varepsilon\)。这次用单位插入中间一层：
\[
\delta_Y=F(\eta_{GY}):QY=FGY\longrightarrow FGFGY=Q^2Y.
\]
其自然性由 \(\eta\) 的自然性经 \(F\) 得到。两条余单位律可以直接核验：
\[
\begin{aligned}
\varepsilon_{QY}\circ\delta_Y
&=\varepsilon_{FGY}\circ F(\eta_{GY})=1_{FGY},\\
Q(\varepsilon_Y)\circ\delta_Y
&=F\bigl(G(\varepsilon_Y)\circ\eta_{GY}\bigr)=1_{FGY}.
\end{aligned}
\]
第一行使用 \(GY\) 处的三角恒等式，第二行使用另一个三角恒等式再施加 \(F\)。对于余结合律，将单位的自然性用于 \(\eta_{GY}:GY\rightarrow GFGY\)，得到
\[
GF(\eta_{GY})\circ\eta_{GY}
=\eta_{GFGY}\circ\eta_{GY}.
\]
施加 \(F\)，两边分别是 \(Q(\delta_Y)\circ\delta_Y\) 与 \(\delta_{QY}\circ\delta_Y\)，故余结合律也成立。这些核验与在相反范畴中应用\cref{b4:prop:adjunction-monad}得到的对偶证明一致。

\ProseParagraphBreak
单子记录可以复合的运算，而由伴随得到它时，另一范畴的全部信息未必都被记录。\Cref{b4:thm:A-em-adjunction}构造单子代数的自由与遗忘伴随；\cref{b4:thm:A-beck-monadicity}再判断原范畴何时能从这些代数恢复。它要求右伴随保守，即右伴随把一个态射送成同构时，该态射本身也必须是同构。另一个条件涉及分裂余等化子：它在余等化图上带有满足特定恒等式的截面与收缩映射，准确条件见\cref{b4:def:A-split-coequalizer}。判据要求每对经右伴随后具有这种分裂图的平行态射，在原范畴中都有余等化子，并被右伴随保持；不能把它理解为伴随自动保持所有余等化子。这是另外的构造与判据，并非伴随本身已经保证的结论。
